\newpage
\phantomsection
\label{prf:sign-preservation}

\begin{remark*}[Return]
\hyperref[prop:sign-preservation]{Return to Theorem}
\end{remark*}

\begin{theorem*}[Sign Preservation]
Let $f : E \to \mathbb{R}$ be continuous at $x_0 \in E$.
If $f(x_0) > 0$, then there exists $\delta > 0$ such that
$f(x) > 0$ for all $x \in E$ with $|x - x_0| < \delta$.
Analogously, if $f(x_0) < 0$, then there exists $\delta > 0$ such that
$f(x) < 0$ for all $x \in E$ with $|x - x_0| < \delta$.
\end{theorem*}

\begin{proof}
\textbf{Professional Standard Proof.}~\\
Without loss of generality, assume $f(x_0) \neq 0$. Set $\varepsilon_0 := |f(x_0)| > 0$. Since $f$ is continuous at $x_0$, there exists $\delta > 0$ such that for all $x \in E$ with $|x - x_0| < \delta$, one has $|f(x) - f(x_0)| < \varepsilon_0 = |f(x_0)|$. This yields $f(x_0) - |f(x_0)| < f(x) < f(x_0) + |f(x_0)|$ for all such $x$. If $f(x_0) > 0$, then $|f(x_0)| = f(x_0)$, so $f(x) > f(x_0) - f(x_0) = 0$. If $f(x_0) < 0$, then $|f(x_0)| = -f(x_0)$, so $f(x) < f(x_0) + (-f(x_0)) = 0$. The same $\delta$ works for both cases.
\end{proof}

\begin{proof}
\textbf{Detailed Learning Proof.}~\\

\textbf{Step 1.} Set up the continuity framework. Since $f(x_0) \neq 0$, define $\varepsilon_0 := |f(x_0)| > 0$. This choice works uniformly for both positive and negative values of $f(x_0)$.

\textbf{Step 2.} Apply continuity at $x_0$. Since $f$ is continuous at $x_0$, for the chosen $\varepsilon_0 = |f(x_0)|$ there exists $\delta > 0$ such that for all $x \in E$,
\[
  |x - x_0| < \delta \;\Rightarrow\; |f(x) - f(x_0)| < |f(x_0)|.
\]

\textbf{Step 3.} Convert the absolute value inequality. The condition $|f(x) - f(x_0)| < |f(x_0)|$ is equivalent to
\[
  -|f(x_0)| < f(x) - f(x_0) < |f(x_0)|,
\]
which rearranges to
\[
  f(x_0) - |f(x_0)| < f(x) < f(x_0) + |f(x_0)|.
\]

\textbf{Step 4.} Handle the case $f(x_0) > 0$. When $f(x_0) > 0$, one has $|f(x_0)| = f(x_0)$. The lower bound from Step 3 gives $f(x) > f(x_0) - f(x_0) = 0$, establishing that $f(x) > 0$ in the $\delta$-neighborhood.

\textbf{Step 5.} Handle the case $f(x_0) < 0$. When $f(x_0) < 0$, one has $|f(x_0)| = -f(x_0)$. The upper bound from Step 3 gives $f(x) < f(x_0) + (-f(x_0)) = 0$, establishing that $f(x) < 0$ in the $\delta$-neighborhood.

\textbf{Step 6.} Conclude. The same $\delta$ obtained in Step 2 serves both cases, completing the proof.
\end{proof}

\begin{remark*}[Proof structure]
This is a direct proof using the definition of continuity. The key insight is choosing $\varepsilon_0 = |f(x_0)|$, which works uniformly for both positive and negative values of $f(x_0)$. The resulting inequality $f(x_0) - |f(x_0)| < f(x) < f(x_0) + |f(x_0)|$ naturally separates into the two cases by substituting the appropriate expression for $|f(x_0)|$. This avoids duplicating the $\delta$ construction.
\end{remark*}

\begin{remark*}[Dependencies]~\\
\begin{itemize}
  \item \hyperref[def:continuous-at-point]{Definition of continuity at a point}
\end{itemize}
\end{remark*}
